\documentclass[instructions.tex]{subfiles}
\begin{document}
 
\chapter{L'injure que le péché fait à Dieu.}
 

Il n'est point d'homme qui osât commettre un seul péché, s'il connaissait l'injure qu'il fait à Dieu. Par le péché, on se révolte contre Dieu; on l'offense avec impudence, avec ingratitude, avec mépris; on l'afflige, on le déshonore. Auriez-vous cru qu'un seul péché mortel renfermât tant de traits de malignité ?

\primo Le péché est une désobéissance et une rébellion. Dieu commande à ses créatures en maître souverain. Il se fait obéir par les astres, par les vents et par les tempêtes; il se fait obéir par les animaux, qui sont soumis à l'homme, parce que Dieu le veut; il se fait obéir par ses anges, qui tremblent à ses volontés adorables; il n'y a que l'homme qui se révolte. Dieu lui commande par sa loi, par les lumières de sa conscience, par son Église; et l'homme dit en son coeur: Non, je n'obéirai point\footnote{Dixisti: Non serviam. Jer.2.}.

\secundo Il se révolte avec l'impudence la plus hardie. Y a-t-il un sujet qui ose outrager son roi en sa présence? Mais l'homme a bien moins de respect pour son Dieu: il l'offense en sa présence et sous ses yeux. Dieu le porte et le soutient entre ses bras\footnote{Portans omnia verbo virtutis suae. Hebr.1.3.}; il peut l'anéantir et le perdre à tout moment: n'importe, l'homme a la hardiesse de s'élever contre lui.

Un homme qu'on tiendrait sur un précipice, aurait-il la témérité d'insulter celui qui aurait ainsi sa vie entre ses mains? C'est vous-même qui êtes si téméraire. Dieu vous soutient par un fil de vie sur un abîme de feu: il peut vous précipiter à chaque instant. Quelle est donc votre fureur d'oser l'outrager? Ah! si vous voulez pécher, cherchez du moins un lieu où Dieu ne puisse vous voir ni vous punir. Mais où trouverez-vous un lieu qui ne soit rempli de sa majesté et de sa gloire?

\tertio En l'offensant, vous vous rendez coupable de la plus noire ingratitude. Dieu vous a donné la vie; il vous a donné votre âme, votre esprit, votre liberté; il vous a donné votre corps et tous vos sens, vos richesses et votre crédit; plus que tout cela, il vous a donné son fils; ce fils adorable vous a donné son sang et ses mérites; et ce sont tous ces bienfaits de Dieu que vous employez contre lui-même, en faisant servir votre esprit, vos sens, ses dons, sa puissance à vos iniquités\footnote{Servire me fecisti in peccatis tuis. Is.43.}.

N'avez-vous donc reçu tant de choses de votre Dieu que pour l'outrager? Et pourquoi de ses dons en faites-vous des armes pour lui faire la guerre? Dieu mérite-t-il qu'on le traite de la sorte? Il ne vous a jamais fait de mal. Eh! quand il vous aurait fait du mal, feriez-vous, pour lui déplaire, plus que vous ne faites? Vous traitez un Dieu, qui vous aime, plus mal que vous ne traiteriez votre ennemi.

Offenser un Dieu qui ne vous a jamais offensé, c'est être bien ingrat. Offenser un Dieu qui ne vous a fait que du bien, c'est être un monstre. Mais offenser un Dieu qui ne vous a fait que du bien, et se servir de ses propres dons pour l'outrager, c'est une ingratitude sur laquelle on ne peut trop verser de larmes.

\quarto Voici un autre trait de malice: c'est que le péché ajoute le mépris à l'ingratitude. Estimer Dieu plus que toutes choses, c'est un devoir; mais l'estimer moins qu'un plaisir honteux, le chasser de votre coeur pour un vil intérêt, pour une vengeance, etc., quel mépris plus injurieux!

Que voulez-vous me donner, et je vous livrerai mon maître? disait Judas aux pontifes. Il le livra pour trente pièces d'argent. Hélas! vous avez trahi Jésus-Christ pour moins de chose, pour contenter une passion, etc.

Vous détestez le procédé des Juifs, qui préférèrent un homme voleur au Fils de Dieu; et vous ne voyez pas que votre conduite est, en un sens, plus détestable. Ce n'est pas un homme, c'est quelque chose de moins; c'est votre sensualité, un point d'honneur, une criminelle attache, que vous préférez à Jésus-Christ. Les Juifs ne connaissaient point ce roi de gloire, et vous le connaissez; vous savez qu'il est votre Sauveur et votre Dieu. Sur quoi donc pouvez-vous excuser le mépris que vous en faites? et pourrez-vous assez pleurer, pour réparer un tel outrage?


\section{Résolutions}

Je déplorerai tous les jours de ma vie le malheur que j'ai eu jusqu'ici d'outrager tant de fois mon Dieu.

\primo Je me rappellerai, surtout dans la tentation, que je ne puis pécher mortellement sans me rendre coupable envers Dieu d'une révolte pleine d'impudence, d'une ingratitude, d'un mépris monstrueux.

\secundo O mon Dieu! je me prosterne aux pieds de Jésus attaché à la croix, et c'est là que j'ose vous demander et attendre miséricorde.

\end{document}
